% 构造函数，d=3，m=3,...,40；不是某个假设空间实测的生长函数。
\begin{tikzpicture}[figure picture]
 \begin{axis}[
  width=108mm,height=78mm,xmin=3,xmax=40,ymin=5,ymax=2e12,ymode=log,
  axis lines=left,xlabel={样本量$m$},ylabel={预测向量数（对数刻度）},
  xtick={3,10,20,30,40},ytick={10,1000,1000000,1000000000,1000000000000},
  log basis y=10,grid=major,
  legend style={at={(.5,1.12)},anchor=south,legend columns=3,column sep=4pt},
  clip=false]
  \addplot[FigureRed,line width=.85pt] coordinates {(3,8) (4,16) (5,32) (6,64) (7,128) (8,256) (9,512) (10,1024) (11,2048) (12,4096) (13,8192) (14,16384) (15,32768) (16,65536) (17,131072) (18,262144) (19,524288) (20,1048576) (21,2097152) (22,4194304) (23,8388608) (24,16777216) (25,33554432) (26,67108864) (27,134217728) (28,268435456) (29,536870912) (30,1073741824) (31,2147483648) (32,4294967296) (33,8589934592) (34,17179869184) (35,34359738368) (36,68719476736) (37,137438953472) (38,274877906944) (39,549755813888) (40,1099511627776)};
  \addlegendentry{$2^m$}
  \addplot[FigureBlue,line width=.85pt,dashed] coordinates {(3,27) (4,64) (5,125) (6,216) (7,343) (8,512) (9,729) (10,1000) (11,1331) (12,1728) (13,2197) (14,2744) (15,3375) (16,4096) (17,4913) (18,5832) (19,6859) (20,8000) (21,9261) (22,10648) (23,12167) (24,13824) (25,15625) (26,17576) (27,19683) (28,21952) (29,24389) (30,27000) (31,29791) (32,32768) (33,35937) (34,39304) (35,42875) (36,46656) (37,50653) (38,54872) (39,59319) (40,64000)};
  \addlegendentry{$m^3$参照}
  \addplot[FigureYellow,line width=.85pt,mark=square*,mark repeat=6,mark size=1.4pt] coordinates {(3,8) (4,15) (5,26) (6,42) (7,64) (8,93) (9,130) (10,176) (11,232) (12,299) (13,378) (14,470) (15,576) (16,697) (17,834) (18,988) (19,1160) (20,1351) (21,1562) (22,1794) (23,2048) (24,2325) (25,2626) (26,2952) (27,3304) (28,3683) (29,4090) (30,4526) (31,4992) (32,5489) (33,6018) (34,6580) (35,7176) (36,7807) (37,8474) (38,9178) (39,9920) (40,10701)};
  \addlegendentry{$\sum_{j=0}^3\binom mj$}
  \draw[FigureMuted,line width=.6pt,{Stealth[length=1.5mm]}-{Stealth[length=1.5mm]}]
   (axis cs:30,4526)--(axis cs:30,1073741824);
 \end{axis}
\end{tikzpicture}%
